\section{Introduction}\label{sec:intro}

Multi-object tracking (MOT) finds every object of interest in a video and gives each one an
identity number (ID) that stays the same over time. It is used in surveillance, driving, and drone
systems \cite{zhu2022visdrone,wu2022uavsurvey}, where results are needed with little delay. Most
MOT systems follow the tracking-by-detection design \cite{adzemovic2026survey}: in each frame, a
detector returns candidate boxes with confidence scores, and a tracker links them into tracks, one
per object \cite{bewley2016sort,zhang2022bytetrack,cao2023ocsort}. The tracker uses a box only if
its score passes a fixed threshold. A two-stage tracker, such as ByteTrack, also uses lower-score
boxes in a second step to continue unmatched tracks \cite{zhang2022bytetrack}; a single-stage
tracker, such as OC-SORT, ignores every box below one threshold \cite{cao2023ocsort}.

In real systems the detector is often replaced, retrained, or compressed, while the tracker keeps
the thresholds that its authors set for another detector. Detector scores are not calibrated across
models and conditions \cite{kuppers2020detcalib}, so a new detector usually gives scores on a
different scale, and nothing in the system notices. The effect can be large: halving the scores of
a published detector reduces two well-known trackers to 0 higher order tracking accuracy (HOTA)
\cite{luiten2021hota}, because no detection passes their thresholds (Sect.~\ref{sec:shift}).
Retuning every tracker for every detector needs labels, which a working system usually does not
have. Adaptive-threshold methods are each built into one tracker
\cite{ma2023adaptivebytetrack,shim2024adaptrack,act2026adaptive}, score calibration needs
labelled data \cite{kuppers2020detcalib}, and adaptive inference methods learn their settings for
one model \cite{han2022dynamic,jiang2018chameleon}.

This paper presents a self-calibrating control layer (\scl{}) that sits between a detector and a
tracker that are both kept fixed (frozen): their code, weights, and stored settings are not
changed. We call the tracker that \scl{} serves the \emph{host}; it gives \scl{} four of its
published settings, the host contract (Sect.~\ref{sec:contract}). \scl{} checks how many of the
higher scores of recent frames form a clearly confident group, using only the scores, so it needs
no labels. If most are confident, the host keeps its scores and, unless they would reject the
confident group, its thresholds, and the layer only adds a step that continues tracks. If many are uncertain, \scl{} chooses which candidates reach the host, gives
them new scores, and sets the host's thresholds for the frame.

With halved detector scores, \scl{} restored two trackers from 0 to about 66 HOTA; with other
detectors, it raised the multiple object tracking accuracy (MOTA) of two-stage trackers by up to
24.7 points and cut their identity switches by 43\% to 56\%, while well-matched strong trackers
stayed within 0.06 HOTA points of their own accuracy.

The main contributions are:
\begin{enumerate}
\item \scl{}, a training-free and causal layer that keeps the tracker's scores and thresholds when
      the scores are clean, lowering a threshold only if it would reject the confident group, and
      takes control only when they are noisy, with a host contract that lets one fixed
      configuration serve single-stage and two-stage trackers without changing them
      (Sect.~\ref{sec:method}).
\item A development path from always-on banding (\aob{}), an earlier design that always
      intervened, in which each step is compared with the step it was built on
      (Sect.~\ref{sec:development}).
\item An evaluation of one frozen configuration with paired bootstrap intervals, comparisons with
      simple adaptive thresholds and the published ACT method, an analysis of each part of the
      layer (Sect.~\ref{sec:results}), and a cost of about 3 ms per frame on an ordinary processor
      (central processing unit, CPU) (Sect.~\ref{sec:cost}).
\end{enumerate}

Section~\ref{sec:related} covers the data, metrics, and related work;
Sections~\ref{sec:method} and~\ref{sec:development} present \scl{} and its development;
Section~\ref{sec:results} gives the results; and Sections~\ref{sec:discussion}
and~\ref{sec:conclusion} discuss and conclude.
